\subsubsection{Continuous Bag-of-Words (CBOW)}\label{sec:continuous-bag-of-words-cbow}

The \definition{Continuous Bag-of-Words (CBOW)} architecture is the first of the two Word2Vec models proposed by Mikolov et al. (2013).\cite{mikolov2013efficient} Its goal is straightforward: \hl{given the surrounding words (the context), predict the missing (central) word.} Unlike Bengio's Neural Language Model, CBOW removes the hidden layer and greatly simplifies the network, allowing it to be trained on billions of words while still learning high-quality word embeddings.

\highspace
Suppose our sentence is:
\begin{center}
    ``\emph{The quick brown fox jumps over the dog.}''
\end{center}

\noindent
And we choose a context window of size 2. Our training example becomes:
\begin{center}
    \emph{Input:} \{``the'', ``brown'', ``jumps'', ``over''\} \\
    \emph{Output:} ``fox''
\end{center}

\noindent
Thus, instead of learning:
\begin{equation*}
    P\left(w_t \, | \, w_{t-2}, w_{t-1}\right)
\end{equation*}
As Bengio's model did, CBOW learns:
\begin{equation*}
    P\left(w_t \, | \, w_{t-c}, \dots w_{t-1}, w_{t+1}, \dots, w_{t+c}\right)
\end{equation*}
Where $w_t$ is the central word, and $c$ is the \textbf{context radius} $\frac{n}{2}$, where $n$ is the \textbf{total number of context words}. In our example, the context words are ``the'', ``brown'', ``jumps'', and ``over'' (two previous and two future words) and the target word is ``fox''. Notice one important improvement over Bengio's model: Bengio used \textbf{previous words}, while CBOW uses \textbf{both previous and future words}, giving more information to predict the target.

\begin{figure}[!htp]
    \centering
    \tikzsetnextfilename{cbow-prediction-context-window}
    \tikzwidth{\textwidth}
    \begin{tikzpicture}[
            >=Latex,
            font=\small,
            word/.style={
                draw,
                rounded corners,
                minimum width=1.25cm,
                minimum height=0.55cm,
                align=center,
                fill=blue!8
            },
            emb/.style={
                draw,
                rounded corners,
                minimum width=1.45cm,
                minimum height=0.55cm,
                align=center,
                fill=green!10
            },
            proj/.style={
                draw,
                rounded corners,
                minimum width=1.35cm,
                minimum height=0.8cm,
                align=center,
                fill=purple!10
            },
            outnode/.style={
                draw,
                rounded corners,
                minimum width=1.45cm,
                minimum height=0.65cm,
                align=center,
                fill=orange!15
            },
            title/.style={font=\bfseries\large, align=center},
            subtitle/.style={font=\small, align=center},
            note/.style={font=\small, align=center, text width=6cm}
        ]

            \node[title] at (0,3.1) {CBOW};
            \node[subtitle] at (0,2.7) {Continuous Bag-of-Words};

            \node[word] (c1) at (-3.8,1.6) {$w_{t-2}$};
            \node[word] (c2) at (-3.8,0.7) {$w_{t-1}$};
            \node[word] (c3) at (-3.8,-0.7) {$w_{t+1}$};
            \node[word] (c4) at (-3.8,-1.6) {$w_{t+2}$};

            \node[emb] (e1) at (-1.9,1.6) {$C(w_{t-2})$};
            \node[emb] (e2) at (-1.9,0.7) {$C(w_{t-1})$};
            \node[emb] (e3) at (-1.9,-0.7) {$C(w_{t+1})$};
            \node[emb] (e4) at (-1.9,-1.6) {$C(w_{t+2})$};

            \node[proj] (avg) at (0.3,0) {
                $\bar{\mathbf v}$\\[-1mm]
                \scriptsize average
            };

            \node[outnode] (soft1) at (2.8,0) {
                Softmax\\[-1mm]
                \scriptsize over $|V|$
            };

            \node[word] (target1) at (5.0,0) {$w_t$};

            \draw[->] (c1) -- (e1);
            \draw[->] (c2) -- (e2);
            \draw[->] (c3) -- (e3);
            \draw[->] (c4) -- (e4);

            \draw[->] (e1) -- (avg);
            \draw[->] (e2) -- (avg);
            \draw[->] (e3) -- (avg);
            \draw[->] (e4) -- (avg);

            \draw[->] (avg) -- (soft1);
            \draw[->] (soft1) -- (target1);

            \node[subtitle] at (-3.8,-2.3) {Input\\context words};
            \node[subtitle] at (-1.9,-2.3) {Embedding\\lookup};
            \node[subtitle] at (0.3,-2.3) {Projection};
            \node[subtitle] at (2.8,-2.3) {Output};
            \node[subtitle] at (5.0,-2.3) {Target};

            \node[note] at (0.4,-3.35) {
            Given the surrounding context words,\\
            CBOW predicts the center word.
            };

            \node[note] at (0.4,-4.35) {
            $\displaystyle
            \bar{\mathbf v}
            =
            \frac{1}{C}
            \sum_{i=1}^{C}
            C(w_i)
            $
            };

        \end{tikzpicture}
    \caption{CBOW architecture in Word2Vec.\cite{mikolov2013efficient}}
    \label{fig:cbow-architecture-repeated}
\end{figure}

\begin{flushleft}
    \textcolor{Green3}{\faIcon{compress} \textbf{Context Aggregation}}
\end{flushleft}
The key idea of CBOW is how it combines the context words. Suppose our vocabulary is represented by the embedding matrix:
\begin{equation*}
    C \in \mathbb{R}^{\left|V\right| \times m}
\end{equation*}
Where $\left|V\right|$ is the vocabulary size and $m$ is the embedding dimension. Each context word $w_i$ performs an embedding lookup in $C$ to retrieve its corresponding vector:
\begin{equation*}
    C\left(w_i\right) \in \mathbb{R}^{m}
\end{equation*}
Unlike Bengio's model, \textbf{these vectors are not concatenated}. Instead, CBOW computes their average:
\begin{equation}
    \bar{\mathbf{v}} = \mathbf{h} = \frac{1}{C} \cdot \sum_{i=1}^{C} C\left(w_i\right)
\end{equation}
Where $C$ is the number of context words. This average vector $\bar{\mathbf{v}}$ is the representation of the \textbf{entire context}.

\newpage

\noindent
\textcolor{Green3}{\faIcon{question-circle} \textbf{Why average?}} Suppose the context is ``\emph{coffee}'', ``\emph{morning}'', ``\emph{drink}''. Their embeddings might encode different semantic aspects:
\begin{itemize}
    \item \emph{coffee} $\rightarrow$ \emph{beverage, caffeine, hot}
    \item \emph{morning} $\rightarrow$ \emph{time, day, early}
    \item \emph{drink} $\rightarrow$ \emph{consume, liquid, beverage}
\end{itemize}
Averaging them produces a vector roughly representing ``\emph{something people drink in the morning}''. This vector is then used to predict the missing word.

\highspace
\textcolor{Green3}{\faIcon{book} \textbf{Output Prediction.}} The averaged vector is fed directly into the output layer, which \hl{computes the probability of every word in the vocabulary.} The predicted word is:
\begin{equation*}
    \hat{w} = \argmax_{w} P(w \, | \, \text{context})
\end{equation*}
During training, the correct target word is known, so the prediction error is computed and the embeddings are updated through backpropagation.

\highspace
\textcolor{Green3}{\faIcon{question-circle} \textbf{Which Embeddings are updated?}} One interesting property of CBOW is that \hl{only the embeddings of the context words are updated.} Suppose the training example is:
\begin{center}
    \emph{Input:} \{``the'', ``brown'', ``jumps'', ``over''\} \\
    \emph{Target:} ``fox''
\end{center}

\noindent
The model retrieves the embeddings of ``the'', ``brown'', ``jumps'', and ``over'' from the embedding matrix $C$. These four embeddings are averaged to produce the prediction. After computing the prediction error, \textbf{only these input (context) embeddings are modified} to better predict \emph{fox}. The target word itself is represented in the output weight matrix (often denoted $D'$ or $C'$), which is updated separately during training.

\highspace
\begin{flushleft}
    \textcolor{Green3}{\faIcon{tachometer-alt} \textbf{Computational Complexity}}
\end{flushleft}
Because CBOW removes the hidden layer, its computational cost is \textbf{dramatically lower than Bengio's NNLM}. For each training example, the complexity is approximately:
\begin{equation}\label{eq:cbow-complexity}
    O\left(n \cdot m + m \cdot \log\left|V\right|\right)
\end{equation}
Where:
\begin{itemize}
    \item $n$ is the number of context words;
    \item $m$ is the embedding dimension;
    \item $\left|V\right|$ is the vocabulary size.
\end{itemize}
Compared to Bengio's complexity:
\begin{equation*}
    O\left(n \cdot m + n \cdot m \cdot h + h \cdot \left|V\right|\right)
\end{equation*}
The expensive \textbf{hidden-layer computations disappear}, making Word2Vec roughly \hl{three orders of magnitude faster} to train on large corpora.